\section{The earlier ordinary free-profile decoder}
\label{sec:sbm-proof}
\label{sbm:section}

We work under the ordinary-SBM assumptions of Theorem~\ref{thm:sbm}.
In particular, $K\ge2$ is fixed, $n_a\ge c_*n$, the block matrix is
symmetric and uniformly full rank at scale $p$, and
$p/2\le P_{ab}\le p\le1/4$. Write $d=np$, $C_a=\{i:z_i=a\}$,
$M_0=\Omega=ZPZ^\top$, and
\[
 \Pi_0=Z(Z^\top Z)^{-1}Z^\top,
 \qquad p_{a*}=(P_{a,z_j})_{j=1}^n.
\]
The diagonal of $M_0$ is retained in these comparison profiles;
the observed graph has no self-loops. The resulting bounded diagonal
correction is included below.

\subsection{A finite decoder using only the retained information}
\label{sbm:algorithm}

Let $D=\diag(A\one)$, and retain $K$ normalized adjacency eigenpairs
$D^{-1/2}AD^{-1/2}U=U\Lambda$, including the Perron direction and
selecting the remaining signal eigenvalues by absolute value. On graphs
where the encoder is defined, put
\[
 X=D^{-1/2}U,\qquad
 Q=X(X^\top X)^{-1/2},\qquad
 Y=DX\Lambda(X^\top X)^{-1/2},\qquad
 R=YQ^\top.
\]
These are exact algebraic operations. The generalized eigenrelation gives
$Y=AQ$ and hence $R=A\Pi$, where $\Pi=QQ^\top$. The Perron direction
implies $\Pi\one=\one$ and $R\one=A\one$. In particular, the entire
decoder below is evaluable from $(D,U,\Lambda)$ without accessing an
individual adjacency entry after compression. The exceptional-graph
convention in the theorem makes the encoder and decoder total.

Fix $0<c_0<1/16$, let $\bar d=n^{-1}\sum_iD_{ii}$, and set
\[
 \varepsilon=c_0\max(1,\bar d)/n,\qquad
 \mathcal I=[\varepsilon,1-\varepsilon].
\]
All clipping in this section is coordinatewise to $\mathcal I$.
For $v\in\R^n$ and $\mu\in\mathcal I^n$, define
\[
 \ell(v;\mu)
 =\sum_j\{v_j\log\mu_j+(1-v_j)\log(1-\mu_j)\},
 \qquad \ell_i(\mu)=\ell(R_{i*};\mu).
\]
The entries of $R$ may be signed. The score is nevertheless finite and
has the retained-factor expression
\begin{equation}
 \ell_i(\mu)=Y_iQ^\top\operatorname{logit}(\mu)
             +\sum_j\log(1-\mu_j).
 \label{sbm:score}
\end{equation}
For $s$ fitted profiles $\mathcal C=(\mu_1,\ldots,\mu_s)$ and simplex
weights $\pi$, use
\[
 \mathcal L(\mathcal C,\pi)
 =\sum_i\log\sum_{a=1}^s\pi_a e^{\ell_i(\mu_a)}.
\]
An E-step and its subsequent M-step are
\begin{align}
 \gamma_{ia}
 &=\frac{\pi_a e^{\ell_i(\mu_a)}}
         {\sum_b\pi_b e^{\ell_i(\mu_b)}},
 &N_a&=\sum_i\gamma_{ia},\notag\\
 \mu_{aj}^{\rm new}
 &=\operatorname{clip}_{\mathcal I}
       \left(\frac{\sum_i\gamma_{ia}R_{ij}}{N_a}\right),
 &\pi_a^{\rm new}&=N_a/n.
 \label{sbm:em}
\end{align}
When $N_a=0$, the old profile is retained and the new weight is zero.
The inner fitting subroutine performs at most $100$ exact E/M updates.
It may stop earlier when its labels are unchanged and its objective
increment is at most $10^{-12}(1+|\mathcal L_{\rm new}|)$. Ties in
all maximizations are settled deterministically. The proof uses only
finiteness and monotonicity of these inner updates.

The candidate attached to vertex $j$ is
$\nu_j=\operatorname{clip}_{\mathcal I}(R_{j*})$. Start from the clipped
global mean and fit one component. For $s=2,\ldots,K$, append a candidate
maximizing
\[
 \sum_i\left(\ell_i(\nu_j)-\max_{a<s}\ell_i(\mu_a)\right)_+,
\]
excluding the vertex identifiers selected in previous growing stages,
and fit all $s$ components from uniform weights. This gives $K$ profiles.
Next perform at most $T=\lceil\log n\rceil$ replacement rounds. For each
deletion $a$ from the round's incumbent, put
\[
 h_i^{(-a)}=\max_{b\ne a}\ell_i(\mu_b)
\]
and choose a candidate maximizing
$\sum_i(\ell_i(\nu_j)-h_i^{(-a)})_+$ over all vertices. Replace profile
$a$ by this candidate, initialize all $K$ weights uniformly, and run
the inner fit. After trying every deletion, retain the fit with largest
$\mathcal L$, including the unchanged incumbent among the choices.
A round without strict improvement terminates replacement.

Finally, evaluate the selected fit's responsibilities, perform one
mandatory M-step in \eqref{sbm:em}, and return
\begin{equation}
 \widehat z_i\in\argmax_a
 \{\ell_i(\mu_a^{\rm new})+\log\pi_a^{\rm new}\}.
 \label{sbm:decode}
\end{equation}
The reference to an E/M update here is ordered: the E-step is evaluated
at the selected fit, its M-step follows, and the decoder uses the new
profiles and weights.

\subsection{A common box and deterministic reconstruction bounds}
\label{sbm:box-section}

Throughout the deterministic arguments, intersect the base event
$\mathcal E_0$ in \eqref{eq:base-event} with the row-concentration
events of Lemma~\ref{lem:normalized-rows}. Its constants are fixed,
independent of $n$. In addition to the degree bounds and normalized
signal gap, its consequences used here are
\begin{equation}
 \norm{A-M_0}_{\op}\le C\sqrt d,\qquad
 \norm{U}_{2\to\infty}\le C/\sqrt n,\qquad
 \norm{\Pi-\Pi_0}_{\op}\le C/\sqrt d.
 \label{sbm:common-event}
\end{equation}
The same conclusions hold for the vertex-deleted graphs on their
corresponding good events. Their probability costs are handled in the
probability closure, after the sharp row calculation.

\begin{lemma}[Reconstruction and generated probability box]
\label{sbm:box}
On the full-graph good event, for all sufficiently large $n$,
\begin{align}
 \norm{R-M_0}_F^2&\le Cd,
 &\max_{i,j}|R_{ij}|&\le Cp,
 &\norm{R_i(I-\Pi_0)}_2&\le C\sqrt p,
 \label{sbm:reconstruction-bounds}\\
 \norm{\Pi_0(A-R)\Pi_0}_{\op}&\le C.
 \label{sbm:block-residual}
\end{align}
There are fixed $a_-,a_+,\eta_->0$ such that every profile generated by
the decoder, every clipped true-class mean, and every true comparison
profile belongs to
\begin{equation}
 \mathcal P_n
 =\{\mu:\ a_-p\le\mu_j\le a_+p,\quad
                    \mu_j\le1-\eta_-\text{ for every }j\}.
 \label{sbm:box-def}
\end{equation}
Every true probability is at distance at least $cp$ from both endpoints
of the actual clipping interval $\mathcal I$.
\end{lemma}

\begin{proof}
Write $E=A-M_0$. Since $R=A\Pi$,
\[
 R-M_0=E\Pi+M_0(\Pi-I).
\]
Each term has rank at most $K$, operator norm at most $C\sqrt d$, and
hence Frobenius norm at most $C\sqrt d$, by
\eqref{sbm:common-event} and $\norm{M_0}_{\op}\le Cd$.
For the other row bounds, write $G=X^\top X$. The degree bounds imply
that the eigenvalues of $G$ are between constant multiples of $d^{-1}$.
Thus $Q=XG^{-1/2}$ has row norms at most $C/\sqrt n$, and
\[
 AQ=DQT_0,\qquad T_0=G^{1/2}\Lambda G^{-1/2},\qquad
 \norm{T_0}_{\op}\le C.
\]
It follows that $\norm{(AQ)_i}_2\le Cd/\sqrt n$, which proves the
entrywise bound for $R=(AQ)Q^\top$. Moreover,
\[
 \norm{R_i(I-\Pi_0)}_2
 \le\norm{(AQ)_i}_2\norm{Q^\top(I-\Pi_0)}_{\op}
 \le C\sqrt{d/n}.
\]

Both $\Pi$ and $\Pi_0$ are orthogonal projectors. Consequently,
\[
 \norm{\Pi_0(I-\Pi)\Pi_0}_{\op}
 =\norm{(I-\Pi)\Pi_0}_{\op}^2\le C/d.
\]
Using $M_0=\Pi_0M_0\Pi_0$ gives
\[
 \Pi_0(A-R)\Pi_0
 =M_0\Pi_0(I-\Pi)\Pi_0+\Pi_0E(I-\Pi)\Pi_0.
\]
The two terms have norms at most $Cd/d$ and
$C\sqrt d/\sqrt d$, respectively. This proves
\eqref{sbm:block-residual} without a commutation assumption on $A$
and $\Pi$.

The average degree satisfies
\[
 \left|\bar d-n^{-1}\one^\top M_0\one\right|
 \le\norm{E}_{\op}\le C\sqrt d.
\]
Since $d/2\le n^{-1}\one^\top M_0\one\le d$, eventually
$d/4\le\bar d\le2d$. In particular $\bar d>1$ eventually, so
\begin{equation}
 c_0p/4\le\varepsilon\le2c_0p<p/8.
 \label{sbm:floor}
\end{equation}
Together with $P_{ab}\in[p/2,p]$ and $p\le1/4$, this gives the claimed
interior distance from the clipping endpoints.

Each active raw M-step coordinate is a convex combination of a column
of $R$, even though that column may have negative entries. It therefore
lies in $[-C_Rp,C_Rp]$. Clipping makes it at most
$\max(C_Rp,\varepsilon)\le C_\mu p$. The initial mean and all clipped
candidates obey the same bound; an inactive component retains an earlier
profile. These observations prove inductively that all generated
profiles lie between $c_0p/4$ and $C_\mu p$. If
$p\le1/(2C_\mu)$, their upper bound is $1/2$. Otherwise the lower
bound in \eqref{sbm:floor} gives
$\varepsilon\ge c_0/(8C_\mu)$, providing a fixed upper gap through
clipping. Enlarge the constants to include the true profiles and the
clipped true-class means.
\end{proof}

The box in \eqref{sbm:box-def} describes profiles generated by this
algorithm and its comparison profiles. Uniform statements below concern
this common box; they do not concern arbitrary vectors in the larger
nominal interval $\mathcal I^n$.

\begin{lemma}[Exact monotonicity for signed observations]
\label{sbm:em-monotonicity}
Every exact update \eqref{sbm:em} increases or preserves $\mathcal L$.
\end{lemma}

\begin{proof}
The log-sum-exp variational identity gives
\[
 \sum_{i,a}\gamma_{ia}
 \{\log\pi_a+\ell_i(\mu_a)-\log\gamma_{ia}\}
 \le\mathcal L(\mathcal C,\pi),
\]
with equality at the E-step. For one profile coordinate and fixed
responsibilities, put $a_j=\sum_i\gamma_{ia}R_{ij}$. Its objective is
\[
 a_j\log\mu+(N_a-a_j)\log(1-\mu),
\]
whose derivative has sign equal to $a_j/N_a-\mu$ when $N_a>0$.
Thus its global maximizer on $\mathcal I$ is the clipped weighted mean,
whether or not $a_j$ lies between zero and $N_a$. The weight maximizer
is $N_a/n$. Zero mass leaves the corresponding coordinate objective
constant. Maximization of this variational lower bound proves the claim.
\end{proof}

\subsection{The finite replacement certificate and common alignment}
\label{sbm:initialization-section}

For a $K$-profile list and weights, define the empirical reference loss
\[
 F_R(\mathcal C,\pi)
 =\sum_i\ell(R_i;p_{z_i,*})-\mathcal L(\mathcal C,\pi),
\]
and its population version
\[
 F_0(\mathcal C,\pi)
 =-\sum_i\log\sum_a\pi_a
                   e^{-\KL(p_{z_i,*}\Vert\mu_a)}.
\]
Here a vector Bernoulli divergence means the coordinate sum
\[
 \KL(p_{c*}\Vert\mu)=\sum_j
 \KL(\Ber(P_{c,z_j})\Vert\Ber(\mu_j)).
\]
The first term defining $F_R$ is a comparison constant independent of
the fitted list and weights. Hence every implemented likelihood
comparison is precisely the reverse comparison of $F_R$. The population
loss $F_0$ is nonnegative. Also put
\begin{align*}
 G_R(\mathcal C)
 &=\sum_i\{\ell(R_i;p_{z_i,*})-\max_a\ell_i(\mu_a)\},\\
 G_0(\mathcal C)
 &=\sum_c n_c\min_a\KL(p_{c*}\Vert\mu_a).
\end{align*}
These empirical reference losses are defined for signed $R$.

\begin{lemma}[Uniform comparison of signed and population objectives]
\label{sbm:objective-comparison}
On the full-graph good event, uniformly over center lists in
$\mathcal P_n$ and simplex weights,
\begin{equation}
 |F_R-F_0|+|G_R-G_0|\le Cn\sqrt d.
 \label{sbm:objective-comparison-bound}
\end{equation}
Always $G_R\le F_R$. With uniform weights,
$F_R\le G_R+n\log K$. Every true class supplies a candidate $\nu_j$
with $\norm{\nu_j-p_{c*}}_1\le C\sqrt d$; all these candidates are
available in every replacement trial.
\end{lemma}

\begin{proof}
For a true class $c$ and fitted profile $\mu_a$, the difference
$\ell(v;\mu_a)-\ell(v;p_{c*})$ is affine in $v$, with coefficient
$\operatorname{logit}(\mu_a)-\operatorname{logit}(p_{c*})$ bounded
uniformly on $\mathcal P_n$. The maximum and log-sum-exp operations are
1-Lipschitz for changes bounded in maximum norm. Therefore each
comparison in \eqref{sbm:objective-comparison-bound} is at most
$C\sum_i\norm{R_i-p_{z_i,*}}_1$. By
\eqref{sbm:reconstruction-bounds}, this sum is at most
$n\norm{R-M_0}_F\le Cn\sqrt d$. The common $\log p$ terms have
canceled in the score contrasts used in this argument.

The inequalities between hard and soft reference losses follow from
\[
 K^{-1}e^{\max_a s_a}\le K^{-1}\sum_a e^{s_a}
 \le e^{\max_a s_a}
\]
for uniform weights, and
$\sum_a\pi_a e^{s_a}\le e^{\max_a s_a}$ for arbitrary weights.
Finally, each class has at least $c_*n$ rows, so one of its rows has
squared error at most $Cd/n$ by the Frobenius bound. Its $\ell_1$
error is at most $C\sqrt d$. Clipping toward an interval containing
the true profile contracts coordinatewise absolute error.
\end{proof}

\begin{lemma}[Finite replacement certificate]
\label{sbm:replacement}
At the selected endpoint of the replacement stage,
\begin{equation}
 F_R\le \xi_n nd,\qquad
 \xi_n=C\{(1-1/K)^{\lceil\log n\rceil}+d^{-1/2}\}\longrightarrow0.
 \label{sbm:certificate}
\end{equation}
The same conclusion holds if replacement stops because a complete
round does not strictly improve the incumbent.
\end{lemma}

\begin{proof}
Assign each true profile a nearest fitted center in population
divergence. A true class contributes at least $G_0/K$. If a fitted
center is unused by this assignment, replacing that center by this
true profile removes its contribution without increasing any other
class's loss. If every center is used, the assignment between the
$K$ true classes and $K$ fitted centers is a bijection; replacing the
center of the highest-loss class has the same property. Consequently
some ideal replacement has population hard loss at most
$(1-1/K)G_0$.

The derivative of $\KL(p_{c*}\Vert\mu)$ in a parameter coordinate is
$(\mu_j-P_{c,z_j})/[\mu_j(1-\mu_j)]$, uniformly bounded on the common
box. Replacing the ideal profile by the available candidate from
Lemma~\ref{sbm:objective-comparison} costs at most $Cn\sqrt d$.
Combining the objective comparisons gives a candidate list $\mathcal C'$
with
\[
 G_R(\mathcal C')\le(1-1/K)F_R(\mathcal C,\pi)+Cn\sqrt d.
\]

For a fixed deletion $a$, the empirical hard loss after inserting
$\nu_j$ is exactly
\[
 \sum_i\{\ell(R_i;p_{z_i,*})-h_i^{(-a)}
                         -(\ell_i(\nu_j)-h_i^{(-a)})_+\}.
\]
Thus the implemented maximum positive gain minimizes this loss among
candidates. All deletions are tried. Uniform trial weights add at most
$n\log K$ to the hard loss, and the subsequent exact inner fit cannot
increase $F_R$ by Lemma~\ref{sbm:em-monotonicity}. Including the
incumbent among the choices proves the recurrence
\begin{equation}
 F_{R,t+1}\le(1-1/K)F_{R,t}+Cn\sqrt d+n\log K,
 \qquad F_{R,t+1}\le F_{R,t}.
 \label{sbm:replacement-recurrence}
\end{equation}

Every generated profile lies in the common box, so each population
divergence is at most $Cd$. Hence $F_{0,0}\le Cnd$ for arbitrary
simplex weights, and the same bound holds for $F_{R,0}$ by
\eqref{sbm:objective-comparison-bound}. No preliminary accuracy guarantee
for the growing stage is needed: its generated list supplies this finite
starting bound. Iterating
\eqref{sbm:replacement-recurrence} for $\lceil\log n\rceil$ rounds
gives \eqref{sbm:certificate}, after enlarging its constant. If a
complete round gives no strict improvement, its chosen loss is the
incumbent's loss. The first inequality in
\eqref{sbm:replacement-recurrence} then immediately yields
$F_R\le K(Cn\sqrt d+n\log K)$, which has the same conclusion.
\end{proof}

\begin{lemma}[Common matching and refined responsibilities]
\label{sbm:responsibilities}
The selected fit admits a single permutation after which
\begin{equation}
 \max_a\norm{\mu_a-p_{a*}}_1=o(d),\qquad
 -\log\min_a\pi_a=o(d).
 \label{sbm:consistent-fit}
\end{equation}
Its next E-step satisfies
\begin{equation}
 \frac1n\sum_{i,a}|\gamma_{ia}-Z_{ia}|\le C/d.
 \label{sbm:responsibility-radius}
\end{equation}
\end{lemma}

\begin{proof}
The preceding lemmas give $F_0\le\zeta_nnd$ with $\zeta_n\to0$.
Since $F_0\ge G_0$, each true class has a fitted center whose
divergence is $o(d)$. On the common box, Bernoulli entropy curvature
and Cauchy--Schwarz imply
\[
 \KL(p_{a*}\Vert\mu)\ge
 \frac{\norm{p_{a*}-\mu}_2^2}{Cp}
 \ge\frac{\norm{p_{a*}-\mu}_1^2}{Cd}.
\]
Thus each true profile has a fitted center at $o(d)$ in $\ell_1$.
Uniform full rank and positive class sizes give
\[
 \min_{a\ne b}\norm{p_{a*}-p_{b*}}_2^2\ge cnp^2,
 \qquad
 \min_{a\ne b}\norm{p_{a*}-p_{b*}}_1\ge cd.
\]
The nearby centers must therefore be distinct. There are exactly $K$
of them, proving one common matching and the first assertion of
\eqref{sbm:consistent-fit}.

For a row of true class $a$, each wrong-center population divergence
is at least $c_1d$ eventually, by this matching and the common-box
Lipschitz bound. Its population soft loss is at least
$-\log(\pi_a+e^{-c_1d})$. It is the same for all $n_a$ rows of that
class, and every population soft loss is nonnegative. Consequently
\[
 \pi_a+e^{-c_1d}\ge e^{-\zeta_nnd/n_a},\qquad
 \pi_a\ge\tfrac12e^{-C\zeta_nd}
\]
for all sufficiently large $n$. This proves the weight assertion.

Fix a sufficiently small constant $t_0>0$. If
$\norm{R_i-p_{a*}}_1\le t_0d$ and $z_i=a$, the population separation,
\eqref{sbm:consistent-fit}, and bounded contrast coefficients give
a fitted correct-class margin at least $c_2d$. The wrong responsibility
mass on this row is at most $(K-1)e^{-c_2d}$. Every remaining row has
squared error at least $t_0^2d^2/n$. The Frobenius bound limits their
number to $Cn/d$. Summing their responsibility masses, which are at
most one, proves \eqref{sbm:responsibility-radius}.
\end{proof}

\subsection{Fitted profiles change the relevant score by a bounded amount}
\label{sbm:fitted-section}

Define the centered block fluctuation
\begin{equation}
 \mathcal B
 =\max_{a,c}
 \left|\frac{\one_{C_a}^\top(A-\E A)\one_{C_c}}
                   {\sqrt{n_an_c}}\right|.
 \label{sbm:block-fluctuation}
\end{equation}

\begin{lemma}[Block moments and deletion]
\label{sbm:block-moments}
For every fixed $t\ge0$, $\E e^{t\mathcal B}\le C_t$. For a tested
vertex $i$, let $\mathcal B_i$ be the analogous fluctuation calculated
from $B=A_{-i,-i}$ and the deleted class sizes. Then
$\E e^{t\mathcal B_i}\le C_t$, and on $D_{ii}\le Cd$,
\begin{equation}
 \mathcal B\le C(1+\mathcal B_i).
 \label{sbm:fluctuation-deletion}
\end{equation}
\end{lemma}

\begin{proof}
Each normalized block sum in \eqref{sbm:block-fluctuation} is a sum
of independent centered upper-triangular Bernoulli entries with
coefficients of magnitude at most $C/n$. Its sum of coefficient-squared
edge probabilities is at most $Cp$. The elementary centered Bernoulli
transform bound therefore gives, for each fixed real $t$,
\[
 \log\E\exp\left\{
 t\frac{\one_{C_a}^\top(A-\E A)\one_{C_c}}
         {\sqrt{n_an_c}}\right\}
 \le C t^2 e^{C|t|/n}p.
\]
For example, this follows by applying
$e^u-1-u\le u^2e^{|u|}/2$ to every independent edge and multiplying
their transforms. Taking both signs and summing over the fixed number
of block pairs bounds $\E e^{t\mathcal B}$.

Deleting one vertex preserves the positive class-size bounds, so the
same argument applies to $\mathcal B_i$. The unnormalized difference
between a full and deleted block sum consists of at most twice the
tested row's centered edge sum. Its absolute value is at most
$C(D_{ii}+d)$. Dividing by $\sqrt{n_an_c}\ge cn$ gives $O(p)$.
The factor converting a deleted normalized block sum to the full
denominator is bounded by a constant. This proves
\eqref{sbm:fluctuation-deletion}.
\end{proof}

For each class define the raw true-class mean and its clipped version
\[
 \bar r_a=n_a^{-1}\sum_{l\in C_a}R_l,
 \qquad \bar\mu_a=\operatorname{clip}_{\mathcal I}(\bar r_a).
\]
These are comparison quantities only. For $a\ne b$, also define
\begin{equation}
 w_j^{ab}=\log\frac{P_{a,z_j}(1-P_{b,z_j})}
                         {P_{b,z_j}(1-P_{a,z_j})},
 \qquad
 r_j^{ab}=\log\frac{1-P_{a,z_j}}{1-P_{b,z_j}}.
 \label{sbm:true-weights}
\end{equation}

\begin{lemma}[Uniform fitted-score comparison]
\label{sbm:fitted-score}
On the full-graph good event, every M-step from responsibilities
satisfying \eqref{sbm:responsibility-radius}, after their common
matching, has weights bounded below and obeys
\begin{equation}
 \left|
 \ell_i(\mu_a)-\ell_i(\mu_b)+\log\frac{\pi_a}{\pi_b}
 -\left\{R_iw^{ab}+\sum_jr_j^{ab}\right\}
 \right|
 \le C(1+\mathcal B)
 \label{sbm:fitted-score-bound}
\end{equation}
uniformly over $i,a,b$ and all such responsibility matrices.
\end{lemma}

\begin{proof}
Write $e=n^{-1}\sum_{l,a}|\gamma_{la}-Z_{la}|\le C/d$.
Then $N_a=n_a+O(ne)$, so $N_a\ge cn$ eventually and the updated
weights $N_a/n$ are bounded below. By
\eqref{sbm:reconstruction-bounds}, $\max_l\norm{R_l}_1\le Cd$.
The difference of the raw responsibility-weighted mean and $\bar r_a$
can be written as a numerator perturbation and a denominator
perturbation. Their combined $\ell_1$ norm is at most
\[
 \frac{Cd}{N_a}\sum_l|\gamma_{la}-Z_{la}|
 +\frac{Cd}{N_a}|N_a-n_a|
 \le Ced\le C.
\]
Clipping is 1-Lipschitz, whence
\begin{equation}
 \norm{\mu_a-\bar\mu_a}_1\le C.
 \label{sbm:actual-perfect-mean}
\end{equation}

Jensen's inequality and the Frobenius bound give
\[
 \norm{\bar r_a-p_{a*}}_2^2
 \le n_a^{-1}\sum_{l\in C_a}\norm{R_l-p_{a*}}_2^2\le Cp.
\]
Clipping toward an interval containing $p_{a*}$ preserves this bound.
The interior distance from Lemma~\ref{sbm:box} also gives the scalar
inequality
\[
 |\operatorname{clip}_{\mathcal I}(x)-x|
 \le\frac{(x-u)^2}{cp}
\]
for each true coordinate $u$. Thus
$\norm{\bar\mu_a-\bar r_a}_1\le C$.

For a true block $C_c$, decompose the block sum of
$\bar r_a-p_{a*}$ using $R-A$, $A-\E A$, and $\E A-M_0$.
The first contributes at most
\[
 \sqrt{n_c/n_a}\,
 \norm{\Pi_0(R-A)\Pi_0}_{\op}\le C,
\]
by \eqref{sbm:block-residual}; the second is at most $C\mathcal B$.
The diagonal correction contributes at most $P_{aa}\le p$.
Including the clipping correction proves, for
$\delta_a=\bar\mu_a-p_{a*}$,
\begin{equation}
 \norm{\delta_a}_2^2\le Cp,
 \qquad
 \left|\sum_{j\in C_c}\delta_{aj}\right|
 \le C(1+\mathcal B).
 \label{sbm:mean-errors}
\end{equation}

For any profile in the common box, the score derivatives are
\[
 \frac{\partial\ell_i}{\partial\mu_j}
 =\frac{R_{ij}-\mu_j}{\mu_j(1-\mu_j)},\qquad
 \frac{\partial^2\ell_i}{\partial\mu_j^2}
 =-\frac{R_{ij}}{\mu_j^2}
  -\frac{1-R_{ij}}{(1-\mu_j)^2}.
\]
Their absolute values are at most $C$ and $C/p$, respectively.
Therefore \eqref{sbm:actual-perfect-mean} changes the score by $O(1)$,
and the second-order remainder in a Taylor expansion at $p_{a*}$ is
at most $C\norm{\delta_a}_2^2/p\le C$.

Split the first-order gradient at $p_{a*}$ into its projection onto
$\col(Z)$ and its orthogonal complement. Its block-constant
projection has bounded coordinates; paired with $\delta_a$, it is
bounded by \eqref{sbm:mean-errors}. The diagonal multiplier
$[P_{a,z_j}(1-P_{a,z_j})]^{-1}$ is block constant and commutes with
$\Pi_0$. Hence the complementary gradient has norm at most
\[
 \frac C p\norm{R_i(I-\Pi_0)}_2\le C/\sqrt p.
\]
Its inner product with $\delta_a$, whose norm is at most $C\sqrt p$,
is $O(1)$. We have proved
$|\ell_i(\mu_a)-\ell_i(p_{a*})|\le C(1+\mathcal B)$.
Apply the same calculation to $b$ and include the bounded updated
log-weight ratio. The exact true-profile contrast is the expression
in braces in \eqref{sbm:fitted-score-bound}, proving the result.
\end{proof}

\subsection{Euclidean projected evidence at the oracle tilt}
\label{sbm:tilted-section}

Fix a tested vertex $i$ and labels $a\ne b$. Let $B=A_{-i,-i}$,
$x=A_{-i,i}$, and let $\Pi_B$ be the Euclidean projector onto the
deleted generalized signal space from Lemma~\ref{lem:loo}. Write
$\Pi_{0,B}$ for the true block projector on the remaining vertices.
For this paragraph, $w$ and $r$ denote the restrictions of
\eqref{sbm:true-weights} to these vertices. The exact oracle log
likelihood ratio is
\[
 L_{i,ab}=b_{i,ab}+x^\top w,
 \qquad b_{i,ab}=\sum_{j\ne i}r_j^{ab}.
\]
Let $I_{i,ab}$ be its exact Bernoulli Chernoff information in
\eqref{eq:sbm-information} and $t_*$
its optimizing parameter, using the transform under class $a$ with
factor $e^{-t_*L_{i,ab}}$. Under the corresponding optimizing tilt,
the row coordinates are independent with block-constant probabilities
\begin{equation}
 q_j=\frac{P_{a,z_j}^{1-t_*}P_{b,z_j}^{t_*}}
 {P_{a,z_j}^{1-t_*}P_{b,z_j}^{t_*}
 +(1-P_{a,z_j})^{1-t_*}(1-P_{b,z_j})^{t_*}}.
 \label{sbm:tilted-probability}
\end{equation}
These probabilities are comparable to $p$. The log-odds are bounded,
and uniform full rank gives $\sum_{j\ne i}w_j^2\asymp n$.
For all tilt parameters in $[0,1]$, the log transform therefore has
second derivative between $cd$ and $Cd$. Its two endpoint values
are zero. It follows that $I_{i,ab}\asymp d$, the optimizing parameter
stays inside a fixed compact subinterval of $(0,1)$, and its tilted
variance is comparable to $d$. Thus the oracle conditions of
Lemma~\ref{lem:sharp} hold uniformly.

\begin{lemma}[Euclidean tilted cancellation]
\label{sbm:tilted-cancellation}
On the deleted-graph good event, let $\beta=(\Pi_B-I)w$. Then
\begin{equation}
 \max_j|\beta_j|\le C,
 \qquad |q^\top\beta|\le C,
 \qquad \sum_jq_j\beta_j^2\le C.
 \label{sbm:tilted-errors}
\end{equation}
Consequently, conditional on $B$, for every real $u$,
\begin{equation}
 \PP_a\{b_{i,ab}+x^\top\Pi_Bw\le u\mid B\}
 \le C\frac{e^{-I_{i,ab}+t_*u}}{\sqrt d}.
 \label{sbm:conditional-tail}
\end{equation}
\end{lemma}

\begin{proof}
The deleted Euclidean basis has row norms at most $C/\sqrt n$, so
bounded entries of $w$ imply bounded entries of $\Pi_Bw$.
Let $\delta=\norm{(I-\Pi_B)\Pi_{0,B}}_{\op}\le C/\sqrt d$,
as follows from \eqref{eq:euclidean-angle} for the deleted graph.
Both $q$ and $w$ belong to $\col(\Pi_{0,B})$. Orthogonality gives
\[
 q^\top\beta
 =-\ip{(I-\Pi_B)q}{(I-\Pi_B)w}.
\]
Since $\norm q_2\le Cp\sqrt n$ and $\norm w_2\le C\sqrt n$,
its absolute value is at most $Cnp\delta^2\le C$.
Also $\norm\beta_2^2\le\delta^2\norm w_2^2\le Cn/d$, so
$\sum_jq_j\beta_j^2\le Cp\norm\beta_2^2\le C$.
This proves \eqref{sbm:tilted-errors}.

The conditional transform identity is
\[
 \E_a e^{-t_*(b_{i,ab}+x^\top\Pi_Bw)}
 =e^{-I_{i,ab}}\E_{q}e^{-t_*x^\top\beta}.
\]
Independence, $\log(1+y)\le y$, and
$e^v-1-v\le v^2e^{|v|}/2$ bound the final log transform by
$C|q^\top\beta|+C\sum_jq_j\beta_j^2\le C$.
Under tilting by this estimated score, the row is still a product of
Bernoulli variables. Its probabilities and complements are comparable
to their oracle-tilted counterparts because the coefficient errors
are bounded. Its score variance is of order $d$: in the oracle-tilted
quadratic form,
\[
 \sum_jq_j(1-q_j)(w_j+\beta_j)^2
 \ge\tfrac12\sum_jq_j(1-q_j)w_j^2
       -\sum_jq_j\beta_j^2\ge cd
\]
eventually, and the corresponding upper bound is $Cd$.
The bounded summands have total third absolute centered moment $O(d)$.
The bounds \eqref{sbm:tilted-errors} give all three premises in
\eqref{eq:tilt-error}, with zero additive intercept error and fixed
constants. Lemma~\ref{lem:sharp} now gives
\eqref{sbm:conditional-tail}.
A threshold $u$ only contributes the displayed factor $e^{t_*u}$;
it does not change this variance or the interval-concentration bound.
\end{proof}

\begin{proposition}[Sharp error bound for the finite fitted decoder]
\label{sbm:good-event-risk}
Align the selected fit by the common matching from
Lemma~\ref{sbm:responsibilities}. Let $\mathcal E_i$ be the intersection
of the full-graph and vertex-deleted good events, including the tested
degree bound, used in Lemmas~\ref{lem:normalized-rows} and~\ref{lem:loo}.
For a vertex with true class $a$ and any $b\ne a$, the final decoder
\eqref{sbm:decode} satisfies
\begin{equation}
 \PP_z\{\widehat z_i\text{ is aligned to }b,\ \mathcal E_i\}
 \le C\frac{e^{-I_{i,ab}}}{\sqrt d}.
 \label{sbm:good-event-bound}
\end{equation}
In particular, summing over competitors and averaging over vertices
bounds the good-event expected misclassification loss by
\begin{equation}
 \frac Cn\sum_i\sum_{b\ne z_i}
              \frac{e^{-I_{i,z_i b}}}{\sqrt d}.
 \label{sbm:good-average}
\end{equation}
\end{proposition}

\begin{proof}
The finite replacement certificate and its mandatory final E/M update
give the premise of Lemma~\ref{sbm:fitted-score}. A wrong assignment
from $a$ to $b$, including an unfavorable score tie, implies on the
full-graph good event that
\[
 R_iw^{ab}+\sum_jr_j^{ab}\le C(1+\mathcal B).
\]
The full-to-deleted Euclidean identity \eqref{eq:loo-euclidean} applies
to the fixed bounded vector $w^{ab}$. Together with
\eqref{sbm:fluctuation-deletion}, it gives the deterministic event
containment
\begin{equation}
 \{\widehat z_i\text{ is aligned to }b,\ \mathcal E_i\}
 \subseteq
 \{b_{i,ab}+x^\top\Pi_Bw_{-i}^{ab}
                \le C(1+\mathcal B_i),\ B\text{ good}\}.
 \label{sbm:event-containment}
\end{equation}
The nonedge intercept changes by $r_i^{ab}=O(p)$ when the tested
diagonal is removed, which is absorbed into the constant.

Use the full-graph event only for this deterministic containment.
The right-hand event does not condition the tested row on full-graph
eigenvectors, fitted parameters, or a degree event. Conditional on
$B$, the incident row has its original independent Bernoulli law.
Applying \eqref{sbm:conditional-tail} with
$u=C(1+\mathcal B_i)$, integrating over good $B$, and using
Lemma~\ref{sbm:block-moments} gives
\[
 \PP_z\{\widehat z_i\text{ is aligned to }b,\ \mathcal E_i\}
 \le C\frac{e^{-I_{i,ab}}}{\sqrt d}
       \E e^{C\mathcal B_i}
 \le C'\frac{e^{-I_{i,ab}}}{\sqrt d}.
\]
This proves \eqref{sbm:good-event-bound}. The common permutation is
the same for every vertex, so the loss minimized over permutations
is at most the loss under this matching. Summing the pairwise bounds
proves \eqref{sbm:good-average}.
\end{proof}

The sharp Bernoulli oracle comparison in Lemma~\ref{lem:sharp}
identifies the terms in \eqref{sbm:good-average} with local oracle
errors up to fixed constants, including the $d^{-1/2}$ factor
\cite{zhouli}. The probability closure pays the complement of these
good events at that same oracle scale and completes
Theorem~\ref{thm:sbm}.

\paragraph{Extension to general positive block profiles.}
The deterministic and conditional arguments also apply to a fixed
positive full-rank block class
\[
 c_-p\le P_{ab}\le c_+p,\qquad
 \max_{a,b}P_{ab}\le1-\epsilon_0,
\]
with positive comparable class sizes and $d\to\infty$; $p$ need not tend
to zero. In particular $p\le p_{\max}:=(1-\epsilon_0)/c_-$.
The average degree is eventually between $c_-d/2$ and $2c_+d$.
Choose
\[
 0<c_0<\min\{1/16,\ c_-/(4c_+),\ \epsilon_0/2\}.
\]
Then $\varepsilon\le2c_0c_+p<c_-p/2$, giving a distance at least
$c_-p/2$ between every true probability and the lower clipping endpoint.
For $n\ge2$ the elementary graph bound $\bar d\le n-1$ also gives
$\varepsilon\le c_0<\epsilon_0/2$. Hence the distance to the upper
endpoint is at least $\epsilon_0/2\ge
\{\epsilon_0/(2p_{\max})\}p$. This proves the required interior
condition for dense as well as sparse profiles.

The generated common-box proof, all likelihood derivatives, and all
block-fluctuation moments now hold with fixed constants depending on
$c_-,c_+,\epsilon_0$. The exact oracle-tilted probabilities remain
between the two candidate probabilities. Both these means and the exact
log-odds vectors are block constant, so the Euclidean tilted calculation
uses no sparse approximation. All arguments above therefore apply on
the proved good event to this general class. Proposition~\ref{prop:sbm-margin}
pays that event's complement under an explicit information--degree
margin; Corollary~\ref{cor:known} supplies the sparse equal-exposure
case where the margin is automatic.
